\section{Conclusion}\label{sec:conclusion}

This work provides the first systematic evaluation of confidence signals for test-time scaling in generative information extraction.
By comparing five signals from three families---structural, surface, and model-internal---across NER and RE tasks of varying difficulty, we disentangle two objectives that prior work has conflated: quality estimation (ranking instances by expected performance) and sample selection (choosing the best output among $N$ candidates).
Our experiments on SciERC, CoNLL-2003, and WNUT-17 with Qwen3-8B and LLaMA-3.1-8B reveal that these objectives require fundamentally different signal properties.

Three findings anchor our analysis.
First, consistency-based signals are effective quality estimators: structural signals (Voting $\rho{=}0.38$, SJ $\rho{=}0.36$) lead rank correlation on complex schemas, while Exact Match achieves the highest binary discrimination (AUROC${=}0.75$) even on hard tasks---revealing that ranking and classification favor different signal families.
Second, the \textit{correlation-selection gap} shows that strong instance-level correlation does not yield selection gains---no consistency signal surpasses greedy decoding on NER best-of-$N$, and only log-probability helps on RE (+2.3~pp)---because instance-level signals cannot discriminate among candidates with similar consistency profiles.
Third, signal rankings depend on task difficulty: exact match leads on easy tasks, Soft Jaccard and voting lead on hard tasks, and all signals converge when model capability is insufficient, establishing practical guidance for signal choice.

Ablations confirm the gap's robustness: it persists at 4B scale and under free-form decoding without grammar constraints (\S\ref{sec:experiments}).

Several directions remain open.
Trained verifiers or reward models could close the selection gap by learning per-sample quality discrimination that unsupervised consistency signals lack.
Cross-task transfer of confidence signals beyond IE---to structured prediction tasks such as semantic parsing and code generation---would test the generality of our findings.
Integration with conformal prediction could convert continuous confidence scores into calibrated coverage guarantees.
Finally, scaling analysis with larger models (70B+) and more diverse IE tasks would clarify whether increased model capability narrows or widens the correlation-selection gap.
